\documentclass[11pt]{article}

\usepackage[margin=1in]{geometry}
\usepackage{amsmath,amssymb,amsthm,mathtools}
\usepackage{booktabs}
\usepackage{enumitem}

\newtheorem{proposition}{Proposition}
\newtheorem{remark}{Remark}
\newtheorem{assumption}{Assumption}

\newcommand{\dd}{\mathrm{d}}
\newcommand{\E}{\mathbb{E}}
\newcommand{\Abar}{\bar A}
\newcommand{\alphabar}{\bar\alpha}
\newcommand{\ibar}{\bar\imath}
\newcommand{\DelA}{\Delta A}
\newcommand{\Delalpha}{\Delta\alpha}

\title{Two-Capital Economy with Adjustment Costs and No Random Capital Shocks}
\author{}
\date{}

\begin{document}
\maketitle

\section{Production}

There are two reproducible capital stocks, dirty capital $K^d_t$ and green capital $K^g_t$.
The two stocks produce a single final good according to linear technologies,
\begin{equation}
    Y_t = A_d K^d_t + A_g K^g_t .
\end{equation}
Investment in capital $j\in\{d,g\}$ is denoted by $I^j_t$, and the investment rate is
\begin{equation}
    i^j_t \equiv \frac{I^j_t}{K^j_t}.
\end{equation}
Consumption is output net of investment:
\begin{equation}
    C_t = (A_d-i^d_t)K^d_t + (A_g-i^g_t)K^g_t .
\end{equation}
The planner maximizes discounted log utility,
\begin{equation}
    \max_{\{i^d_t,i^g_t\}_{t\geq 0}}
    \int_0^\infty e^{-\delta t}\delta \log C_t\,\dd t .
\end{equation}
The factor $\delta$ multiplying flow utility is a normalization. It does not change the optimal controls.

\section{Deterministic Adjustment-Cost Capital Dynamics}

In this note, there are adjustment costs but no random terms in capital accumulation. Capital evolves according to
\begin{equation}
    \frac{\dd K^j_t}{K^j_t}
    = \phi_j(i^j_t)\,\dd t,
    \qquad j\in\{d,g\},
\end{equation}
where
\begin{equation}
    \phi_j(i)
    \equiv
    \alpha_j+\Gamma_j\log(1+\theta_j i).
\end{equation}
The natural domain restriction is
\begin{equation}
    i^j_t> -\frac{1}{\theta_j},
\end{equation}
because $\log(1+\theta_j i^j_t)$ must be well-defined.

The parameters $\alpha_j$, $\Gamma_j$, and $\theta_j$ have the following interpretations. The parameter $\alpha_j$ is the baseline drift of capital $j$, $\Gamma_j$ controls the productivity of investment in generating installed capital, and $\theta_j$ controls the curvature of the adjustment technology.

\section{State Reduction}

Define total capital and the green-capital share by
\begin{equation}
    K_t \equiv K^d_t+K^g_t,
    \qquad
    Z_t \equiv \frac{K^g_t}{K^d_t+K^g_t}.
\end{equation}
Then
\begin{equation}
    K^d_t=(1-Z_t)K_t,
    \qquad
    K^g_t=Z_tK_t.
\end{equation}
Define scaled consumption as
\begin{equation}
    c_t \equiv \frac{C_t}{K_t}.
\end{equation}
Using the resource constraint,
\begin{equation}
    c_t
    =
    (1-Z_t)(A_d-i^d_t)+Z_t(A_g-i^g_t).
\end{equation}
It is useful to define average productivity and average investment by
\begin{equation}
    \Abar(Z)\equiv (1-Z)A_d+ZA_g,
    \qquad
    \ibar\equiv (1-Z)i^d+Zi^g.
\end{equation}
Then the resource constraint can be written compactly as
\begin{equation}
    c=\Abar(Z)-\ibar.
\end{equation}

\section{Derivation of the $\log K$ and $Z$ Processes}

Since
\begin{equation}
    \dd K_t=\dd K^d_t+\dd K^g_t,
\end{equation}
we have
\begin{align}
    \frac{\dd K_t}{K_t}
    &=
    \frac{K^d_t}{K_t}\frac{\dd K^d_t}{K^d_t}
    +
    \frac{K^g_t}{K_t}\frac{\dd K^g_t}{K^g_t} \\
    &=
    (1-Z_t)\phi_d(i^d_t)\,\dd t
    +
    Z_t\phi_g(i^g_t)\,\dd t.
\end{align}
Because the dynamics are deterministic, $\dd\log K_t=\dd K_t/K_t$. Hence
\begin{equation}
    \boxed{
    \frac{\dd \log K_t}{\dd t}
    =
    (1-Z_t)\phi_d(i^d_t)+Z_t\phi_g(i^g_t).
    }
\end{equation}

Next, using
\begin{equation}
    Z_t=\frac{K^g_t}{K_t},
\end{equation}
we obtain
\begin{align}
    \frac{\dd Z_t}{\dd t}
    &=
    \frac{\dot K^g_t K_t-K^g_t\dot K_t}{K_t^2} \\
    &=
    Z_t\phi_g(i^g_t)
    -
    Z_t\left[(1-Z_t)\phi_d(i^d_t)+Z_t\phi_g(i^g_t)\right] \\
    &=
    Z_t(1-Z_t)\left[\phi_g(i^g_t)-\phi_d(i^d_t)\right].
\end{align}
Therefore
\begin{equation}
    \boxed{
    \frac{\dd Z_t}{\dd t}
    =
    Z_t(1-Z_t)\left[\phi_g(i^g_t)-\phi_d(i^d_t)\right].
    }
\end{equation}
The green share rises when green capital grows faster than dirty capital:
\begin{equation}
    \phi_g(i^g_t)>\phi_d(i^d_t)
    \quad\Longrightarrow\quad
    \dot Z_t>0.
\end{equation}

\section{HJB Equation}

Let the value function be written in the reduced state variables as
\begin{equation}
    V(\log K,Z).
\end{equation}
The HJB equation is
\begin{align}
    0
    =
    \max_{i^d,i^g}
    \Bigg\{
    &\delta\log C
    -\delta V(\log K,Z)
    +V_{\log K}(\log K,Z)
    \left[(1-Z)\phi_d(i^d)+Z\phi_g(i^g)\right]
    \nonumber\\
    &+V_Z(\log K,Z)
    Z(1-Z)\left[\phi_g(i^g)-\phi_d(i^d)\right]
    \Bigg\}.
\end{align}
Since
\begin{equation}
    C=Kc,
    \qquad
    \log C=\log K+\log c,
\end{equation}
and the model is homogeneous of degree one in total capital, conjecture
\begin{equation}
    V(\log K,Z)=\log K+v(Z).
\end{equation}
Then
\begin{equation}
    V_{\log K}=1,
    \qquad
    V_Z=v'(Z).
\end{equation}
Substituting into the HJB and canceling the $\log K$ terms gives the one-dimensional deterministic HJB:
\begin{equation}
    \boxed{
    \begin{aligned}
    0
    =
    \max_{i^d,i^g}
    \Big\{
    &\delta\log c
    -\delta v(Z)
    +(1-Z)\phi_d(i^d)+Z\phi_g(i^g)
    \\
    &+Z(1-Z)\left[\phi_g(i^g)-\phi_d(i^d)\right]v'(Z)
    \Big\}.
    \end{aligned}
    }
\end{equation}

\section{Marginal Values and an Equivalent HJB}

Define the marginal-value elasticities
\begin{equation}
    q_d(Z)
    \equiv
    K\frac{\partial V}{\partial K^d},
    \qquad
    q_g(Z)
    \equiv
    K\frac{\partial V}{\partial K^g}.
\end{equation}
Using
\begin{equation}
    \frac{\partial Z}{\partial K^d}=-\frac{Z}{K},
    \qquad
    \frac{\partial Z}{\partial K^g}=\frac{1-Z}{K},
\end{equation}
we obtain
\begin{equation}
    \boxed{
    q_d(Z)=1-Zv'(Z),
    \qquad
    q_g(Z)=1+(1-Z)v'(Z).
    }
\end{equation}
These objects satisfy the Euler identity
\begin{equation}
    \boxed{
    (1-Z)q_d(Z)+Zq_g(Z)=1.
    }
\end{equation}
Using $q_d$ and $q_g$, the HJB can be written as
\begin{equation}
    \boxed{
    0
    =
    \max_{i^d,i^g}
    \left\{
    \delta\log c
    -\delta v(Z)
    +(1-Z)q_d(Z)\phi_d(i^d)
    +Zq_g(Z)\phi_g(i^g)
    \right\}.
    }
\end{equation}
This form is often the cleanest for deriving first-order conditions.

\section{First-Order Conditions}

The control variables are the investment rates $i^d$ and $i^g$. The first-order condition is not taken with respect to $K^d$ or $K^g$; rather, the marginal values $q_d$ and $q_g$ enter because investment changes the growth rate of the corresponding capital stock.

From scaled consumption,
\begin{equation}
    c=(1-Z)(A_d-i^d)+Z(A_g-i^g),
\end{equation}
we have
\begin{equation}
    \frac{\partial c}{\partial i^d}=-(1-Z),
    \qquad
    \frac{\partial c}{\partial i^g}=-Z.
\end{equation}
Also,
\begin{equation}
    \phi'_j(i^j)=\frac{\Gamma_j\theta_j}{1+\theta_j i^j}.
\end{equation}
The first-order conditions are
\begin{equation}
    -\frac{\delta(1-Z)}{c}
    +(1-Z)q_d(Z)\phi'_d(i^d)=0,
\end{equation}
\begin{equation}
    -\frac{\delta Z}{c}
    +Zq_g(Z)\phi'_g(i^g)=0.
\end{equation}
After canceling the common factors $(1-Z)$ and $Z$, the FOCs become
\begin{equation}
    \boxed{
    \frac{\delta}{c}
    =
    q_d(Z)\frac{\Gamma_d\theta_d}{1+\theta_d i^d},
    \qquad
    \frac{\delta}{c}
    =
    q_g(Z)\frac{\Gamma_g\theta_g}{1+\theta_g i^g}.
    }
\end{equation}
Equivalently,
\begin{equation}
    \boxed{
    i^d
    =
    \frac{\Gamma_d c q_d(Z)}{\delta}-\frac{1}{\theta_d},
    \qquad
    i^g
    =
    \frac{\Gamma_g c q_g(Z)}{\delta}-\frac{1}{\theta_g}.
    }
\end{equation}
The admissibility conditions require
\begin{equation}
    q_d(Z)>0,
    \qquad
    q_g(Z)>0,
\end{equation}
or equivalently
\begin{equation}
    -\frac{1}{1-Z}<v'(Z)<\frac{1}{Z}.
\end{equation}

\section{General Adjustment Technologies}

For general $\Gamma_d$, $\Gamma_g$, $\theta_d$, and $\theta_g$, substituting the FOC-implied controls into the resource constraint gives
\begin{equation}
    \boxed{
    c(Z)
    =
    \frac{
    \delta\left[\Abar(Z)+\dfrac{1-Z}{\theta_d}+\dfrac{Z}{\theta_g}\right]
    }{
    \delta+(1-Z)\Gamma_d q_d(Z)+Z\Gamma_g q_g(Z)
    }.
    }
\end{equation}
Therefore, in the fully general case, scaled consumption still depends on the unknown value-function slope $v'(Z)$ through $q_d$ and $q_g$.

\section{Common Adjustment Technology}

Now impose the common adjustment technology used in the calibration:
\begin{equation}
    \Gamma_d=\Gamma_g=\Gamma,
    \qquad
    \theta_d=\theta_g=\theta.
\end{equation}
Then the Euler identity implies an exact analytical aggregation result.

The FOCs imply
\begin{equation}
    1+\theta i^d
    =
    \frac{\Gamma\theta c}{\delta}q_d(Z),
    \qquad
    1+\theta i^g
    =
    \frac{\Gamma\theta c}{\delta}q_g(Z).
\end{equation}
Taking the capital-weighted average gives
\begin{align}
    1+\theta\ibar
    &=
    (1-Z)(1+\theta i^d)+Z(1+\theta i^g) \\
    &=
    \frac{\Gamma\theta c}{\delta}
    \left[(1-Z)q_d+Zq_g\right] \\
    &=
    \frac{\Gamma\theta c}{\delta}.
\end{align}
Combining this equation with $c=\Abar(Z)-\ibar$ gives
\begin{equation}
    \boxed{
    c(Z)
    =
    \frac{\delta\left[1+\theta\Abar(Z)\right]}{\theta(\delta+\Gamma)}.
    }
\end{equation}
The aggregate investment rate is
\begin{equation}
    \boxed{
    \ibar(Z)
    =
    \frac{\Gamma\theta\Abar(Z)-\delta}{\theta(\delta+\Gamma)}.
    }
\end{equation}
Thus aggregate consumption and aggregate investment are available in closed form.

The individual investment rates satisfy
\begin{equation}
    \boxed{
    1+\theta i^d
    =
    \frac{\Gamma\left[1+\theta\Abar(Z)\right]}{\delta+\Gamma}q_d(Z),
    }
\end{equation}
\begin{equation}
    \boxed{
    1+\theta i^g
    =
    \frac{\Gamma\left[1+\theta\Abar(Z)\right]}{\delta+\Gamma}q_g(Z).
    }
\end{equation}
Therefore
\begin{equation}
    \boxed{
    \frac{1+\theta i^g}{1+\theta i^d}
    =
    \frac{q_g(Z)}{q_d(Z)}
    =
    \frac{1+(1-Z)v'(Z)}{1-Zv'(Z)}.
    }
\end{equation}
The investment-rate wedge is
\begin{equation}
    \boxed{
    i^g-i^d
    =
    \frac{\Gamma c(Z)}{\delta}v'(Z).
    }
\end{equation}
Equivalently,
\begin{equation}
    \boxed{
    i^d=\ibar-Z(i^g-i^d),
    \qquad
    i^g=\ibar+(1-Z)(i^g-i^d).
    }
\end{equation}

\section{Value-Function Equation Under Common Adjustment Costs}

All controls are closed form once $v'(Z)$ is known. The remaining problem is to determine $v(Z)$.

Using the FOC,
\begin{equation}
    1+\theta i^j
    =
    \frac{\Gamma\theta c(Z)}{\delta}q_j(Z),
\end{equation}
so
\begin{equation}
    \phi_j(i^j)
    =
    \alpha_j
    +
    \Gamma\log\left(
    \frac{\Gamma\theta c(Z)}{\delta}q_j(Z)
    \right).
\end{equation}
Define
\begin{equation}
    \alphabar(Z)\equiv (1-Z)\alpha_d+Z\alpha_g,
    \qquad
    \Delalpha\equiv \alpha_g-\alpha_d.
\end{equation}
Substituting the optimal controls into the HJB gives the nonlinear first-order ODE
\begin{equation}
    \boxed{
    \begin{aligned}
    0=
    &\delta\log c(Z)
    +\alphabar(Z)
    +\Gamma\log\frac{\Gamma\theta c(Z)}{\delta}
    +\Gamma\left[(1-Z)\log q_d(Z)+Z\log q_g(Z)\right]
    -\delta v(Z)
    \\
    &+Z(1-Z)
    \left[
    \Delalpha
    +\Gamma\log\frac{q_g(Z)}{q_d(Z)}
    \right]v'(Z),
    \end{aligned}
    }
\end{equation}
where
\begin{equation}
    q_d(Z)=1-Zv'(Z),
    \qquad
    q_g(Z)=1+(1-Z)v'(Z).
\end{equation}
This equation is first order because there is no random capital shock. With random capital shocks, a diffusion term in $Z$ would appear and the equation would become second order.

\section{Analytical Solution: What Is Closed Form?}

The general asymmetric model does not usually admit a closed-form solution for $v(Z)$ because the value-function equation is nonlinear in both $v(Z)$ and $v'(Z)$. However, the following objects are analytical under common adjustment costs:
\begin{align}
    c(Z)
    &=
    \frac{\delta\left[1+\theta\Abar(Z)\right]}{\theta(\delta+\Gamma)}, \\
    \ibar(Z)
    &=
    \frac{\Gamma\theta\Abar(Z)-\delta}{\theta(\delta+\Gamma)}, \\
    i^g-i^d
    &=
    \frac{\Gamma c(Z)}{\delta}v'(Z).
\end{align}
Thus, the only unknown object needed for the investment split is the value-function slope $v'(Z)$.

\subsection{Exact Symmetric Benchmark}

If
\begin{equation}
    A_d=A_g=A,
    \qquad
    \alpha_d=\alpha_g=\alpha,
\end{equation}
then the composition of capital is irrelevant and
\begin{equation}
    v'(Z)=0.
\end{equation}
Therefore
\begin{equation}
    q_d(Z)=q_g(Z)=1,
\end{equation}
and
\begin{equation}
    i^d=i^g=i^*.
\end{equation}
The closed-form consumption ratio is
\begin{equation}
    c^*
    =
    \frac{\delta(1+\theta A)}{\theta(\delta+\Gamma)}.
\end{equation}
The closed-form investment rate is
\begin{equation}
    i^*
    =
    \frac{\Gamma\theta A-\delta}{\theta(\delta+\Gamma)}.
\end{equation}
The value function is constant in $Z$ and satisfies
\begin{equation}
    \delta v
    =
    \delta\log c^*
    +\alpha
    +\Gamma\log\frac{\Gamma\theta c^*}{\delta}.
\end{equation}
Hence
\begin{equation}
    \boxed{
    v
    =
    \log c^*
    +\frac{\alpha}{\delta}
    +\frac{\Gamma}{\delta}\log\frac{\Gamma\theta c^*}{\delta}.
    }
\end{equation}

\section{Perturbation Approximation}

Now suppose the two capitals are close to symmetric. Define
\begin{equation}
    \DelA\equiv A_g-A_d,
    \qquad
    \Delalpha\equiv \alpha_g-\alpha_d.
\end{equation}
Treat $\DelA$ and $\Delalpha$ as small. Let $W(A,\alpha)$ denote the symmetric value from the previous section:
\begin{equation}
    W(A,\alpha)
    =
    \log c^*(A)
    +\frac{\alpha}{\delta}
    +\frac{\Gamma}{\delta}\log\frac{\Gamma\theta c^*(A)}{\delta},
\end{equation}
where
\begin{equation}
    c^*(A)=\frac{\delta(1+\theta A)}{\theta(\delta+\Gamma)}.
\end{equation}
A first-order approximation is
\begin{equation}
    \boxed{
    v(Z)
    \approx
    W\left(\Abar(Z),\alphabar(Z)\right).
    }
\end{equation}
Differentiating gives
\begin{equation}
    v'(Z)
    \approx
    W_A\left(\Abar(Z),\alphabar(Z)\right)\DelA
    +
    W_\alpha\left(\Abar(Z),\alphabar(Z)\right)\Delalpha.
\end{equation}
The envelope derivatives are
\begin{equation}
    W_A(A,\alpha)=\frac{1}{c^*(A)},
    \qquad
    W_\alpha(A,\alpha)=\frac{1}{\delta}.
\end{equation}
Since $c(Z)=c^*(\Abar(Z))$, the perturbation approximation is
\begin{equation}
    \boxed{
    v'(Z)
    \approx
    \frac{A_g-A_d}{c(Z)}
    +
    \frac{\alpha_g-\alpha_d}{\delta}.
    }
\end{equation}
The corresponding investment-rate wedge is
\begin{equation}
    \boxed{
    i^g-i^d
    \approx
    \frac{\Gamma}{\delta}(A_g-A_d)
    +
    \frac{\Gamma c(Z)}{\delta^2}(\alpha_g-\alpha_d).
    }
\end{equation}
If $\alpha_d=\alpha_g$, this simplifies to
\begin{equation}
    \boxed{
    i^g-i^d
    \approx
    \frac{\Gamma}{\delta}(A_g-A_d).
    }
\end{equation}
This expression is useful because it shows directly how a green-productivity advantage affects the optimal investment split. A higher $A_g$ raises the marginal value of green capital, which raises $i^g$ relative to $i^d$.

\section{Finite Difference Method for the Deterministic HJB}

The exact value function can be computed by solving the nonlinear first-order HJB on a grid for $Z\in[0,1]$. Since the equation is deterministic and first order, an upwind finite difference scheme is preferable to a centered difference scheme.

\subsection{Step 1: Choose a Grid}

Choose $N+1$ grid points
\begin{equation}
    Z_n=n\Delta Z,
    \qquad
    n=0,1,\ldots,N,
    \qquad
    \Delta Z=\frac{1}{N}.
\end{equation}
Let $v_n$ approximate $v(Z_n)$.

\subsection{Step 2: Set Boundary Values}

At $Z=0$, the economy contains only dirty capital. The boundary value is the one-capital dirty value:
\begin{equation}
    v_0
    =
    \log c_d
    +\frac{\alpha_d}{\delta}
    +\frac{\Gamma}{\delta}\log\frac{\Gamma\theta c_d}{\delta},
\end{equation}
where
\begin{equation}
    c_d=\frac{\delta(1+\theta A_d)}{\theta(\delta+\Gamma)}.
\end{equation}
At $Z=1$, the economy contains only green capital. The boundary value is
\begin{equation}
    v_N
    =
    \log c_g
    +\frac{\alpha_g}{\delta}
    +\frac{\Gamma}{\delta}\log\frac{\Gamma\theta c_g}{\delta},
\end{equation}
where
\begin{equation}
    c_g=\frac{\delta(1+\theta A_g)}{\theta(\delta+\Gamma)}.
\end{equation}

\subsection{Step 3: Initialize the Interior}

A useful initial guess is the perturbation approximation:
\begin{equation}
    v_n^{(0)}=W\left(\Abar(Z_n),\alphabar(Z_n)\right),
    \qquad n=1,\ldots,N-1.
\end{equation}
This provides a smooth starting point close to the exact solution when heterogeneity is moderate.

\subsection{Step 4: Construct One-Sided Derivatives}

For an interior node $n$, define the backward and forward differences:
\begin{equation}
    D^-v_n=\frac{v_n-v_{n-1}}{\Delta Z},
    \qquad
    D^+v_n=\frac{v_{n+1}-v_n}{\Delta Z}.
\end{equation}
For any candidate slope $p$, define
\begin{equation}
    q_d(p;Z_n)=1-Z_np,
    \qquad
    q_g(p;Z_n)=1+(1-Z_n)p.
\end{equation}
Given $p$, compute
\begin{equation}
    c_n=\frac{\delta\left[1+\theta\Abar(Z_n)\right]}{\theta(\delta+\Gamma)},
\end{equation}
\begin{equation}
    i^d_n(p)=\frac{\Gamma c_n q_d(p;Z_n)}{\delta}-\frac{1}{\theta},
    \qquad
    i^g_n(p)=\frac{\Gamma c_n q_g(p;Z_n)}{\delta}-\frac{1}{\theta}.
\end{equation}
Then compute the drift of $Z$:
\begin{equation}
    \mu_Z(p;Z_n)
    =
    Z_n(1-Z_n)
    \left[
    \phi_g(i^g_n(p))-\phi_d(i^d_n(p))
    \right].
\end{equation}

\subsection{Step 5: Use an Upwind Derivative}

The upwind slope is chosen according to the sign of the drift:
\begin{equation}
    p_n=
    \begin{cases}
    D^-v_n, & \text{if } \mu_Z(Z_n)\geq 0,\\
    D^+v_n, & \text{if } \mu_Z(Z_n)<0.
    \end{cases}
\end{equation}
The intuition is simple. If $\mu_Z>0$, the state moves to the right, so the relevant derivative is the backward derivative. If $\mu_Z<0$, the state moves to the left, so the relevant derivative is the forward derivative.

Because $\mu_Z$ depends on the derivative $p_n$, one can choose the upwind direction using the drift from the previous iteration, or solve the nonlinear upwind equations directly with Newton's method.

\subsection{Step 6: Form the Nonlinear Residual}

At each interior node, define the residual
\begin{equation}
    R_n(v)
    =
    \delta\log c_n
    -\delta v_n
    +(1-Z_n)\phi_d(i^d_n)
    +Z_n\phi_g(i^g_n)
    +
    \mu_Z(Z_n) p_n.
\end{equation}
The finite difference approximation to the HJB is
\begin{equation}
    R_n(v)=0,
    \qquad n=1,\ldots,N-1.
\end{equation}
Together with the two boundary conditions, this gives a nonlinear system for
\begin{equation}
    (v_0,v_1,\ldots,v_N).
\end{equation}

\subsection{Step 7: Solve the Nonlinear System}

There are two common approaches.

\paragraph{Newton method.}
Solve
\begin{equation}
    R(v)=0
\end{equation}
using Newton or a damped Newton method. Damping is useful because the residual is nonlinear in $v'$ through $q_d$, $q_g$, and the logarithms.

\paragraph{Policy iteration.}
Given a guess $v^{(m)}$, compute the slope, the implied controls, and the drift. Then solve the linear equation
\begin{equation}
    \delta v^{(m+1)}_n
    =
    \delta\log c_n
    +(1-Z_n)\phi_d(i^{d,(m)}_n)
    +Z_n\phi_g(i^{g,(m)}_n)
    +\mu^{(m)}_{Z,n}p^{(m+1)}_n,
\end{equation}
using the same upwind direction. Iterate until
\begin{equation}
    \max_n |v^{(m+1)}_n-v^{(m)}_n|
\end{equation}
is below a chosen tolerance.

\subsection{Step 8: Recover the Controls}

Once the fixed point $v_n$ is obtained, compute the numerical derivative $p_n$ using the upwind rule. Then compute
\begin{equation}
    q_{d,n}=1-Z_np_n,
    \qquad
    q_{g,n}=1+(1-Z_n)p_n,
\end{equation}
\begin{equation}
    c_n=\frac{\delta\left[1+\theta\Abar(Z_n)\right]}{\theta(\delta+\Gamma)},
\end{equation}
\begin{equation}
    i^d_n=\frac{\Gamma c_n q_{d,n}}{\delta}-\frac{1}{\theta},
    \qquad
    i^g_n=\frac{\Gamma c_n q_{g,n}}{\delta}-\frac{1}{\theta}.
\end{equation}
Finally, aggregate investment is
\begin{equation}
    \ibar_n=(1-Z_n)i^d_n+Z_ni^g_n,
\end{equation}
and the investment wedge is
\begin{equation}
    i^g_n-i^d_n=\frac{\Gamma c_n}{\delta}p_n.
\end{equation}

\section{Summary}

For the deterministic two-capital adjustment-cost model:
\begin{enumerate}[label=(\roman*)]
    \item Total capital evolves according to
    \begin{equation}
        \dot{\log K}=(1-Z)\phi_d(i^d)+Z\phi_g(i^g).
    \end{equation}
    \item The green-capital share evolves according to
    \begin{equation}
        \dot Z=Z(1-Z)[\phi_g(i^g)-\phi_d(i^d)].
    \end{equation}
    \item The reduced value function satisfies
    \begin{equation}
        V(\log K,Z)=\log K+v(Z).
    \end{equation}
    \item The FOCs are
    \begin{equation}
        \frac{\delta}{c}=q_d\phi'_d(i^d),
        \qquad
        \frac{\delta}{c}=q_g\phi'_g(i^g).
    \end{equation}
    \item Under common adjustment costs, $c(Z)$ and $\ibar(Z)$ are exactly closed form.
    \item The investment split is exactly closed form given $v'(Z)$.
    \item The exact asymmetric value function solves a nonlinear first-order ODE.
    \item The first-order perturbation gives
    \begin{equation}
        v'(Z)\approx \frac{A_g-A_d}{c(Z)}+\frac{\alpha_g-\alpha_d}{\delta},
    \end{equation}
    and therefore
    \begin{equation}
        i^g-i^d
        \approx
        \frac{\Gamma}{\delta}(A_g-A_d)
        +
        \frac{\Gamma c(Z)}{\delta^2}(\alpha_g-\alpha_d).
    \end{equation}
    \item The exact solution can be computed with an upwind finite difference method for the first-order deterministic HJB.
\end{enumerate}

\end{document}
